\cleardoublepage
\phantomsection

\addcontentsline{toc}{chapter}{Abstract}
%\chapter*{Abstract}

\begin{center}
    {\Large \textbf{\thesistitle}} \vspace{0.2cm}

    {\large by} \\
    \textsc{\authorname} \vspace{0.2cm}

    Master's dissertation submitted in order to obtain the academic degree of \degree \par\vspace{0.4cm}

    Supervisor: \supervisor \\
    Counsellor: \counsellor \vspace{0.4cm}

    \university \\
    \faculty \\
    \department \\
\end{center}

\vspace{0.6cm}

\noindent\textbf{Abstract} ---
Linear precoding has become a cornerstone of \acs*{MIMO} communications, as it enables spatial multiplexing gains that translate directly into higher spectral efficiency.
These benefits are particularly valuable in multi-beam high-throughput satellite systems, where aggressive frequency reuse across the narrow spot beams turns the forward link into an interference-limited \acs*{MU}-\acs*{MIMO} channel. 
In this environment, precoding is crucial for mitigating inter-beam interference.
In mobile satellite scenarios, however, the combination of long propagation delays and time-varying channels induced by user mobility introduce complications.
More specific, the \acs*{CSI} available during the precoder computation is generally outdated with respect to the channel experienced at the time of transmission.
This work investigates the performance degradation caused by this \acs*{CSI} mismatch and proposes a low-complexity compensation strategy.

\textbf{Part I} reviews the theoretical foundations underlying this work. 
After introducing the \acs*{SU}- and \acs*{MU}-\acs*{MIMO} system models, together with the most relevant information-theoretic concepts, a generalised \acs*{WMMSE} framework for joint optimisation of the precoder and combiner in a \acs*{SU} scenario is presented, and the capacity-achieving precoder is derived.
For the \acs*{MU}-\acs*{MIMO} broadcast channel, we break down common linear precoding strategies and compare there performance.
Although these results are well established, they are revisited here in a unified notation and validated through Monte-Carlo simulations, so that they can serve as a consistent baseline for the satellite-specific analysis that follows.

\textbf{Part II} applies this framework to the forward link of a multi-beam mobile satellite system.
The impact of the \acs*{CSI} feedback delay on the system performance is quantified, revealing a substantial performance loss.
To counter this degradation, a low-complexity prediction strategy based is proposed.
Numerical results show that this scheme recovers a significant fraction of the performance loss under the considered channel model.

\vspace{0.4cm}

\noindent\textbf{Keywords} --- 
\acs*{MIMO}, Precoding Techniques, Mobile Satellite Systems, Outdated \acs*{CSI}
